% --- Archon physics formalization source begin ---
% archon:physics
% archon:covers PhyXMiniProblems/problem_phyx_mini_0716.lean
% archon:source-report reports/phyx_mini/problem_phyx_mini_0716.source.json

\chapter{Physics problem phyx\_mini\_0716}
\label{ch:PhyXMiniProblems_problem_phyx_mini_0716}

\paragraph{Problem source.}
\#\# Physical scenario

Astronomers discover a binary star system with a period of \textbackslash{}(90\textbackslash{},\textbackslash{}mathrm\{days\}\textbackslash{}). Both stars have a mass twice that of the sun.

\#\# Question

How far apart are the two stars?

\#\# Answer choices

- A: A: \$8.3\textbackslash{}times10\textasciicircum{}\{10\}\textbackslash{}

- B: B: \$9.3\textbackslash{}times10\textasciicircum{}\{9\}\textbackslash{}

- C: C: \$9.3\textbackslash{}times10\textasciicircum{}\{10\}\textbackslash{}

- D: D: \$1.03\textbackslash{}times10\textasciicircum{}\{11\}\textbackslash{}

\#\# Figure caption (auxiliary; use the image as primary evidence)

The image depicts a diagram of a binary star system. It includes two stars labeled 1 and 2, each revolving around a central point, which represents the center of mass. The stars are shown in circular orbits with the radius of each orbit labeled as \textbackslash{}( r \textbackslash{}). 

Two arrows labeled \textbackslash{}(\textbackslash{}vec\{F\}\_\{2 \textbackslash{}text\{ on \} 1\}\textbackslash{}) and \textbackslash{}(\textbackslash{}vec\{F\}\_\{1 \textbackslash{}text\{ on \} 2\}\textbackslash{}) represent the gravitational forces between the two stars, directed towards each other. Text indicates that the distance between the stars is \textbackslash{}( d = 2r \textbackslash{}).

Additional text explains: "Both stars revolve around the center of mass in an orbit with radius \textbackslash{}( r \textbackslash{})."

Dashed blue arrows and text emphasize directions and relational distances between elements.

\#\# Dataset metadata

Dynamics; Physical Model Grounding Reasoning, Multi-Formula Reasoning

\paragraph{Recorded answer/context.}
C: C: \$9.3\textbackslash{}times10\textasciicircum{}\{10\}\textbackslash{}

\paragraph{Figure/image path.}
/root/proposal\_for\_physic/science-mango/phyx\_mini\_run/phyx\_data/test\_image/716.png

\paragraph{Formalization target.}
create a compiling Lean file with sorry bodies at `PhyXMiniProblems/problem\_phyx\_mini\_0716.lean`. The Lean declarations must preserve the physical quantities, dimensions or dimensional roles, figure labels, governing-law hypotheses, and final relation expressed by this problem.
Use Mathlib/Physlib names found through LeanExplore where available. If a physics API is missing, introduce faithful local abstractions rather than scalar placeholder aliases.

\begin{theorem}[Physics formalization target]
\label{thm:physics:phyx_mini_0716:target}
\lean{PhyXMiniProblems.ProblemPhyXMini0716.problem_phyx_mini_0716}
\uses{def:physics:phyx-mini-0716:phyxminiproblems-problemphyxmini0716-binarystarsystem, def:physics:phyx-mini-0716:phyxminiproblems-problemphyxmini0716-matchesprimarybinarystarfigure, def:physics:phyx-mini-0716:phyxminiproblems-problemphyxmini0716-matchesbinarystarscenario, def:physics:phyx-mini-0716:phyxminiproblems-problemphyxmini0716-usesstandardgravitationalconstant, def:physics:phyx-mini-0716:phyxminiproblems-problemphyxmini0716-hasphysicalbinarystarparameters, def:physics:phyx-mini-0716:phyxminiproblems-problemphyxmini0716-satisfiesnewtoniancircularorbitlaws, def:physics:phyx-mini-0716:phyxminiproblems-problemphyxmini0716-recordedanswerchoice, def:physics:phyx-mini-0716:phyxminiproblems-problemphyxmini0716-roundstodisplayeddistance, def:physics:phyx-mini-0716:phyxminiproblems-problemphyxmini0716-isuniqueclosestanswer}
The assigned autoformalize agent should translate this physics problem into Lean declarations in the covered file, with theorem and lemma proof bodies written as `by sorry`.
\end{theorem}
\begin{proof}
This is an autoformalization task, not a proof task. Produce statements that can later be proved by the physics prover without weakening the physical model.
\end{proof}

% --- Archon declaration topology begin ---
\section{Declaration topology}
\begin{definition}[gravitational Constant Dimension]
  \label{def:physics:phyx-mini-0716:phyxminiproblems-problemphyxmini0716-gravitationalconstantdimension}
  \lean{PhyXMiniProblems.ProblemPhyXMini0716.gravitationalConstantDimension}
  The physical dimension L\textasciicircum{}3 M\textasciicircum{}-1 T\textasciicircum{}-2 of Newton's gravitational constant
\end{definition}
\begin{proof}
  This is the declared model definition of gravitational Constant Dimension.
\end{proof}
\begin{definition}[force Dimension]
  \label{def:physics:phyx-mini-0716:phyxminiproblems-problemphyxmini0716-forcedimension}
  \lean{PhyXMiniProblems.ProblemPhyXMini0716.forceDimension}
  The physical dimension M L T\textasciicircum{}-2 of force
\end{definition}
\begin{proof}
  This is the declared model definition of force Dimension.
\end{proof}
\begin{definition}[Mass Quantity]
  \label{def:physics:phyx-mini-0716:phyxminiproblems-problemphyxmini0716-massquantity}
  \lean{PhyXMiniProblems.ProblemPhyXMini0716.MassQuantity}
  A nonnegative physical mass
\end{definition}
\begin{proof}
  This is the declared model definition of Mass Quantity.
\end{proof}
\begin{definition}[Duration Quantity]
  \label{def:physics:phyx-mini-0716:phyxminiproblems-problemphyxmini0716-durationquantity}
  \lean{PhyXMiniProblems.ProblemPhyXMini0716.DurationQuantity}
  A nonnegative physical duration
\end{definition}
\begin{proof}
  This is the declared model definition of Duration Quantity.
\end{proof}
\begin{definition}[Length Quantity]
  \label{def:physics:phyx-mini-0716:phyxminiproblems-problemphyxmini0716-lengthquantity}
  \lean{PhyXMiniProblems.ProblemPhyXMini0716.LengthQuantity}
  A nonnegative physical length
\end{definition}
\begin{proof}
  This is the declared model definition of Length Quantity.
\end{proof}
\begin{definition}[Force Magnitude Quantity]
  \label{def:physics:phyx-mini-0716:phyxminiproblems-problemphyxmini0716-forcemagnitudequantity}
  \lean{PhyXMiniProblems.ProblemPhyXMini0716.ForceMagnitudeQuantity}
  \uses{def:physics:phyx-mini-0716:phyxminiproblems-problemphyxmini0716-forcedimension}
  A nonnegative physical force magnitude
\end{definition}
\begin{proof}
  This is the declared model definition of Force Magnitude Quantity.
\end{proof}
\begin{definition}[Gravitational Constant Quantity]
  \label{def:physics:phyx-mini-0716:phyxminiproblems-problemphyxmini0716-gravitationalconstantquantity}
  \lean{PhyXMiniProblems.ProblemPhyXMini0716.GravitationalConstantQuantity}
  \uses{def:physics:phyx-mini-0716:phyxminiproblems-problemphyxmini0716-gravitationalconstantdimension}
  A nonnegative physical value of Newton's gravitational constant
\end{definition}
\begin{proof}
  This is the declared model definition of Gravitational Constant Quantity.
\end{proof}
\begin{definition}[nominal Solar Mass Unit Choices]
  \label{def:physics:phyx-mini-0716:phyxminiproblems-problemphyxmini0716-nominalsolarmassunitchoices}
  \lean{PhyXMiniProblems.ProblemPhyXMini0716.nominalSolarMassUnitChoices}
  Unit choices that differ from SI only by measuring mass in nominal solar masses
\end{definition}
\begin{proof}
  This is the declared model definition of nominal Solar Mass Unit Choices.
\end{proof}
\begin{definition}[day Unit Choices]
  \label{def:physics:phyx-mini-0716:phyxminiproblems-problemphyxmini0716-dayunitchoices}
  \lean{PhyXMiniProblems.ProblemPhyXMini0716.dayUnitChoices}
  Unit choices that differ from SI only by measuring time in 24-hour days
\end{definition}
\begin{proof}
  This is the declared model definition of day Unit Choices.
\end{proof}
\begin{definition}[nonnegative Readout]
  \label{def:physics:phyx-mini-0716:phyxminiproblems-problemphyxmini0716-nonnegativereadout}
  \lean{PhyXMiniProblems.ProblemPhyXMini0716.nonnegativeReadout}
  Read a nonnegative dimensionful quantity in the specified coherent units
\end{definition}
\begin{proof}
  This is the declared model definition of nonnegative Readout.
\end{proof}
\begin{definition}[mass In Kilograms]
  \label{def:physics:phyx-mini-0716:phyxminiproblems-problemphyxmini0716-massinkilograms}
  \lean{PhyXMiniProblems.ProblemPhyXMini0716.massInKilograms}
  \uses{def:physics:phyx-mini-0716:phyxminiproblems-problemphyxmini0716-massquantity, def:physics:phyx-mini-0716:phyxminiproblems-problemphyxmini0716-nonnegativereadout}
  Kilogram readout of a physical mass
\end{definition}
\begin{proof}
  This is the declared model definition of mass In Kilograms.
\end{proof}
\begin{definition}[mass In Nominal Solar Masses]
  \label{def:physics:phyx-mini-0716:phyxminiproblems-problemphyxmini0716-massinnominalsolarmasses}
  \lean{PhyXMiniProblems.ProblemPhyXMini0716.massInNominalSolarMasses}
  \uses{def:physics:phyx-mini-0716:phyxminiproblems-problemphyxmini0716-massquantity, def:physics:phyx-mini-0716:phyxminiproblems-problemphyxmini0716-nominalsolarmassunitchoices, def:physics:phyx-mini-0716:phyxminiproblems-problemphyxmini0716-nonnegativereadout}
  Nominal-solar-mass readout of a physical mass
\end{definition}
\begin{proof}
  This is the declared model definition of mass In Nominal Solar Masses.
\end{proof}
\begin{definition}[duration In Seconds]
  \label{def:physics:phyx-mini-0716:phyxminiproblems-problemphyxmini0716-durationinseconds}
  \lean{PhyXMiniProblems.ProblemPhyXMini0716.durationInSeconds}
  \uses{def:physics:phyx-mini-0716:phyxminiproblems-problemphyxmini0716-durationquantity, def:physics:phyx-mini-0716:phyxminiproblems-problemphyxmini0716-nonnegativereadout}
  Second readout of a physical duration
\end{definition}
\begin{proof}
  This is the declared model definition of duration In Seconds.
\end{proof}
\begin{definition}[duration In Days]
  \label{def:physics:phyx-mini-0716:phyxminiproblems-problemphyxmini0716-durationindays}
  \lean{PhyXMiniProblems.ProblemPhyXMini0716.durationInDays}
  \uses{def:physics:phyx-mini-0716:phyxminiproblems-problemphyxmini0716-durationquantity, def:physics:phyx-mini-0716:phyxminiproblems-problemphyxmini0716-dayunitchoices, def:physics:phyx-mini-0716:phyxminiproblems-problemphyxmini0716-nonnegativereadout}
  24-hour-day readout of a physical duration
\end{definition}
\begin{proof}
  This is the declared model definition of duration In Days.
\end{proof}
\begin{definition}[length In Meters]
  \label{def:physics:phyx-mini-0716:phyxminiproblems-problemphyxmini0716-lengthinmeters}
  \lean{PhyXMiniProblems.ProblemPhyXMini0716.lengthInMeters}
  \uses{def:physics:phyx-mini-0716:phyxminiproblems-problemphyxmini0716-lengthquantity, def:physics:phyx-mini-0716:phyxminiproblems-problemphyxmini0716-nonnegativereadout}
  Metre readout of a physical length
\end{definition}
\begin{proof}
  This is the declared model definition of length In Meters.
\end{proof}
\begin{definition}[force In Newtons]
  \label{def:physics:phyx-mini-0716:phyxminiproblems-problemphyxmini0716-forceinnewtons}
  \lean{PhyXMiniProblems.ProblemPhyXMini0716.forceInNewtons}
  \uses{def:physics:phyx-mini-0716:phyxminiproblems-problemphyxmini0716-forcemagnitudequantity, def:physics:phyx-mini-0716:phyxminiproblems-problemphyxmini0716-nonnegativereadout}
  Newton readout of a physical force magnitude
\end{definition}
\begin{proof}
  This is the declared model definition of force In Newtons.
\end{proof}
\begin{definition}[gravitational Constant In SI]
  \label{def:physics:phyx-mini-0716:phyxminiproblems-problemphyxmini0716-gravitationalconstantinsi}
  \lean{PhyXMiniProblems.ProblemPhyXMini0716.gravitationalConstantInSI}
  \uses{def:physics:phyx-mini-0716:phyxminiproblems-problemphyxmini0716-gravitationalconstantquantity, def:physics:phyx-mini-0716:phyxminiproblems-problemphyxmini0716-nonnegativereadout}
  SI readout of G, in cubic metres per kilogram per second squared
\end{definition}
\begin{proof}
  This is the declared model definition of gravitational Constant In SI.
\end{proof}
\begin{definition}[Star]
  \label{def:physics:phyx-mini-0716:phyxminiproblems-problemphyxmini0716-star}
  \lean{PhyXMiniProblems.ProblemPhyXMini0716.Star}
  The two stars numbered in the supplied diagram
\end{definition}
\begin{proof}
  This is the declared model definition of Star.
\end{proof}
\begin{definition}[Directed Force Label]
  \label{def:physics:phyx-mini-0716:phyxminiproblems-problemphyxmini0716-directedforcelabel}
  \lean{PhyXMiniProblems.ProblemPhyXMini0716.DirectedForceLabel}
  The two directed force labels printed in the diagram
\end{definition}
\begin{proof}
  This is the declared model definition of Directed Force Label.
\end{proof}
\begin{definition}[acting Star]
  \label{def:physics:phyx-mini-0716:phyxminiproblems-problemphyxmini0716-directedforcelabel-actingstar}
  \lean{PhyXMiniProblems.ProblemPhyXMini0716.DirectedForceLabel.actingStar}
  \uses{def:physics:phyx-mini-0716:phyxminiproblems-problemphyxmini0716-star, def:physics:phyx-mini-0716:phyxminiproblems-problemphyxmini0716-directedforcelabel}
  The star on which a labelled force acts
\end{definition}
\begin{proof}
  This is the declared model definition of acting Star.
\end{proof}
\begin{definition}[exerting Star]
  \label{def:physics:phyx-mini-0716:phyxminiproblems-problemphyxmini0716-directedforcelabel-exertingstar}
  \lean{PhyXMiniProblems.ProblemPhyXMini0716.DirectedForceLabel.exertingStar}
  \uses{def:physics:phyx-mini-0716:phyxminiproblems-problemphyxmini0716-star, def:physics:phyx-mini-0716:phyxminiproblems-problemphyxmini0716-directedforcelabel}
  The star exerting a labelled force
\end{definition}
\begin{proof}
  This is the declared model definition of exerting Star.
\end{proof}
\begin{definition}[force Acting On]
  \label{def:physics:phyx-mini-0716:phyxminiproblems-problemphyxmini0716-forceactingon}
  \lean{PhyXMiniProblems.ProblemPhyXMini0716.forceActingOn}
  \uses{def:physics:phyx-mini-0716:phyxminiproblems-problemphyxmini0716-star, def:physics:phyx-mini-0716:phyxminiproblems-problemphyxmini0716-directedforcelabel}
  The force label whose arrow acts on a given star
\end{definition}
\begin{proof}
  This is the declared model definition of force Acting On.
\end{proof}
\begin{definition}[Force Arrow Direction]
  \label{def:physics:phyx-mini-0716:phyxminiproblems-problemphyxmini0716-forcearrowdirection}
  \lean{PhyXMiniProblems.ProblemPhyXMini0716.ForceArrowDirection}
  Directions available for the mutually attractive force arrows
\end{definition}
\begin{proof}
  This is the declared model definition of Force Arrow Direction.
\end{proof}
\begin{definition}[inward Direction]
  \label{def:physics:phyx-mini-0716:phyxminiproblems-problemphyxmini0716-directedforcelabel-inwarddirection}
  \lean{PhyXMiniProblems.ProblemPhyXMini0716.DirectedForceLabel.inwardDirection}
  \uses{def:physics:phyx-mini-0716:phyxminiproblems-problemphyxmini0716-directedforcelabel, def:physics:phyx-mini-0716:phyxminiproblems-problemphyxmini0716-forcearrowdirection}
  Direction dictated by attraction for each printed force label
\end{definition}
\begin{proof}
  This is the declared model definition of inward Direction.
\end{proof}
\begin{definition}[Orbit Direction]
  \label{def:physics:phyx-mini-0716:phyxminiproblems-problemphyxmini0716-orbitdirection}
  \lean{PhyXMiniProblems.ProblemPhyXMini0716.OrbitDirection}
  The sense of revolution indicated by the tangential arrows in the image
\end{definition}
\begin{proof}
  This is the declared model definition of Orbit Direction.
\end{proof}
\begin{definition}[Binary Star Figure]
  \label{def:physics:phyx-mini-0716:phyxminiproblems-problemphyxmini0716-binarystarfigure}
  \lean{PhyXMiniProblems.ProblemPhyXMini0716.BinaryStarFigure}
  \uses{def:physics:phyx-mini-0716:phyxminiproblems-problemphyxmini0716-star, def:physics:phyx-mini-0716:phyxminiproblems-problemphyxmini0716-directedforcelabel, def:physics:phyx-mini-0716:phyxminiproblems-problemphyxmini0716-forcearrowdirection, def:physics:phyx-mini-0716:phyxminiproblems-problemphyxmini0716-orbitdirection}
  Qualitative evidence transcribed from image 716.png. These Boolean fields record which literal marks and labels are shown; physical radii and distances are independent quantities in BinaryStarSystem below
\end{definition}
\begin{proof}
  This is the declared model definition of Binary Star Figure.
\end{proof}
\begin{definition}[Binary Star System]
  \label{def:physics:phyx-mini-0716:phyxminiproblems-problemphyxmini0716-binarystarsystem}
  \lean{PhyXMiniProblems.ProblemPhyXMini0716.BinaryStarSystem}
  \uses{def:physics:phyx-mini-0716:phyxminiproblems-problemphyxmini0716-massquantity, def:physics:phyx-mini-0716:phyxminiproblems-problemphyxmini0716-durationquantity, def:physics:phyx-mini-0716:phyxminiproblems-problemphyxmini0716-lengthquantity, def:physics:phyx-mini-0716:phyxminiproblems-problemphyxmini0716-forcemagnitudequantity, def:physics:phyx-mini-0716:phyxminiproblems-problemphyxmini0716-gravitationalconstantquantity, def:physics:phyx-mini-0716:phyxminiproblems-problemphyxmini0716-star, def:physics:phyx-mini-0716:phyxminiproblems-problemphyxmini0716-directedforcelabel, def:physics:phyx-mini-0716:phyxminiproblems-problemphyxmini0716-binarystarfigure}
  Independent observables of the binary system. In particular, separation is not defined from any answer choice, and no numerical separation appears in this structure
\end{definition}
\begin{proof}
  This is the declared model definition of Binary Star System.
\end{proof}
\begin{definition}[Matches Primary Binary Star Figure]
  \label{def:physics:phyx-mini-0716:phyxminiproblems-problemphyxmini0716-matchesprimarybinarystarfigure}
  \lean{PhyXMiniProblems.ProblemPhyXMini0716.MatchesPrimaryBinaryStarFigure}
  \uses{def:physics:phyx-mini-0716:phyxminiproblems-problemphyxmini0716-lengthinmeters, def:physics:phyx-mini-0716:phyxminiproblems-problemphyxmini0716-binarystarsystem}
  Literal image readout and its associated geometry. The equality d = 2r is printed in the supplied figure. It is geometric input, not the requested numerical separation
\end{definition}
\begin{proof}
  This is the declared model definition of Matches Primary Binary Star Figure.
\end{proof}
\begin{definition}[Matches Binary Star Scenario]
  \label{def:physics:phyx-mini-0716:phyxminiproblems-problemphyxmini0716-matchesbinarystarscenario}
  \lean{PhyXMiniProblems.ProblemPhyXMini0716.MatchesBinaryStarScenario}
  \uses{def:physics:phyx-mini-0716:phyxminiproblems-problemphyxmini0716-massinkilograms, def:physics:phyx-mini-0716:phyxminiproblems-problemphyxmini0716-massinnominalsolarmasses, def:physics:phyx-mini-0716:phyxminiproblems-problemphyxmini0716-durationinseconds, def:physics:phyx-mini-0716:phyxminiproblems-problemphyxmini0716-durationindays, def:physics:phyx-mini-0716:phyxminiproblems-problemphyxmini0716-binarystarsystem}
  The numerical data stated in the problem, with coherent SI conversions made explicit for the later Newtonian calculation. The value of a nominal solar mass is the Physlib value 1.988416 * 10\textasciicircum{}30 kg
\end{definition}
\begin{proof}
  This is the declared model definition of Matches Binary Star Scenario.
\end{proof}
\begin{definition}[Uses Standard Gravitational Constant]
  \label{def:physics:phyx-mini-0716:phyxminiproblems-problemphyxmini0716-usesstandardgravitationalconstant}
  \lean{PhyXMiniProblems.ProblemPhyXMini0716.UsesStandardGravitationalConstant}
  \uses{def:physics:phyx-mini-0716:phyxminiproblems-problemphyxmini0716-gravitationalconstantinsi, def:physics:phyx-mini-0716:phyxminiproblems-problemphyxmini0716-binarystarsystem}
  SI calibration of Newton's gravitational constant used in the calculation
\end{definition}
\begin{proof}
  This is the declared model definition of Uses Standard Gravitational Constant.
\end{proof}
\begin{definition}[Has Physical Binary Star Parameters]
  \label{def:physics:phyx-mini-0716:phyxminiproblems-problemphyxmini0716-hasphysicalbinarystarparameters}
  \lean{PhyXMiniProblems.ProblemPhyXMini0716.HasPhysicalBinaryStarParameters}
  \uses{def:physics:phyx-mini-0716:phyxminiproblems-problemphyxmini0716-massinkilograms, def:physics:phyx-mini-0716:phyxminiproblems-problemphyxmini0716-durationinseconds, def:physics:phyx-mini-0716:phyxminiproblems-problemphyxmini0716-lengthinmeters, def:physics:phyx-mini-0716:phyxminiproblems-problemphyxmini0716-forceinnewtons, def:physics:phyx-mini-0716:phyxminiproblems-problemphyxmini0716-gravitationalconstantinsi, def:physics:phyx-mini-0716:phyxminiproblems-problemphyxmini0716-binarystarsystem}
  Nondegeneracy conditions for the Newtonian circular-orbit model
\end{definition}
\begin{proof}
  This is the declared model definition of Has Physical Binary Star Parameters.
\end{proof}
\begin{definition}[Satisfies Newtonian Circular Orbit Laws]
  \label{def:physics:phyx-mini-0716:phyxminiproblems-problemphyxmini0716-satisfiesnewtoniancircularorbitlaws}
  \lean{PhyXMiniProblems.ProblemPhyXMini0716.SatisfiesNewtonianCircularOrbitLaws}
  \uses{def:physics:phyx-mini-0716:phyxminiproblems-problemphyxmini0716-massinkilograms, def:physics:phyx-mini-0716:phyxminiproblems-problemphyxmini0716-durationinseconds, def:physics:phyx-mini-0716:phyxminiproblems-problemphyxmini0716-lengthinmeters, def:physics:phyx-mini-0716:phyxminiproblems-problemphyxmini0716-forceinnewtons, def:physics:phyx-mini-0716:phyxminiproblems-problemphyxmini0716-gravitationalconstantinsi, def:physics:phyx-mini-0716:phyxminiproblems-problemphyxmini0716-forceactingon, def:physics:phyx-mini-0716:phyxminiproblems-problemphyxmini0716-binarystarsystem}
  Newtonian laws for the idealized equal-mass circular binary: the mutual force has magnitude G m₁ m₂ / d²; the force on each star supplies m ω² r, where ω = 2π/T; and the two labelled forces form an action-reaction pair. These are general governing relations involving the independent observable d; no answer choice or numerical value of d occurs here
\end{definition}
\begin{proof}
  This is the declared model definition of Satisfies Newtonian Circular Orbit Laws.
\end{proof}
\begin{theorem}[binary Star Separation Cubed]
  \label{thm:physics:phyx-mini-0716:phyxminiproblems-problemphyxmini0716-binarystarseparationcubed}
  \lean{PhyXMiniProblems.ProblemPhyXMini0716.binaryStarSeparationCubed}
  \uses{def:physics:phyx-mini-0716:phyxminiproblems-problemphyxmini0716-massinkilograms, def:physics:phyx-mini-0716:phyxminiproblems-problemphyxmini0716-durationinseconds, def:physics:phyx-mini-0716:phyxminiproblems-problemphyxmini0716-lengthinmeters, def:physics:phyx-mini-0716:phyxminiproblems-problemphyxmini0716-gravitationalconstantinsi, def:physics:phyx-mini-0716:phyxminiproblems-problemphyxmini0716-binarystarsystem, def:physics:phyx-mini-0716:phyxminiproblems-problemphyxmini0716-matchesprimarybinarystarfigure, def:physics:phyx-mini-0716:phyxminiproblems-problemphyxmini0716-hasphysicalbinarystarparameters, def:physics:phyx-mini-0716:phyxminiproblems-problemphyxmini0716-satisfiesnewtoniancircularorbitlaws}
  Eliminating the force and orbit radius from the governing laws yields the two-body form of Kepler's third law for the separation d
\end{theorem}
\begin{proof}
  The conclusion follows from the governing relations and figure readouts named in the statement of binary Star Separation Cubed.
\end{proof}
\begin{definition}[Answer Choice]
  \label{def:physics:phyx-mini-0716:phyxminiproblems-problemphyxmini0716-answerchoice}
  \lean{PhyXMiniProblems.ProblemPhyXMini0716.AnswerChoice}
  Labels of the four displayed multiple-choice distances
\end{definition}
\begin{proof}
  This is the declared model definition of Answer Choice.
\end{proof}
\begin{definition}[distance In Meters]
  \label{def:physics:phyx-mini-0716:phyxminiproblems-problemphyxmini0716-answerchoice-distanceinmeters}
  \lean{PhyXMiniProblems.ProblemPhyXMini0716.AnswerChoice.distanceInMeters}
  \uses{def:physics:phyx-mini-0716:phyxminiproblems-problemphyxmini0716-answerchoice}
  Distance in metres printed beside each answer choice
\end{definition}
\begin{proof}
  This is the declared model definition of distance In Meters.
\end{proof}
\begin{definition}[recorded Answer Choice]
  \label{def:physics:phyx-mini-0716:phyxminiproblems-problemphyxmini0716-recordedanswerchoice}
  \lean{PhyXMiniProblems.ProblemPhyXMini0716.recordedAnswerChoice}
  \uses{def:physics:phyx-mini-0716:phyxminiproblems-problemphyxmini0716-answerchoice}
  The answer label recorded by the supplied dataset
\end{definition}
\begin{proof}
  This is the declared model definition of recorded Answer Choice.
\end{proof}
\begin{definition}[Rounds To Displayed Distance]
  \label{def:physics:phyx-mini-0716:phyxminiproblems-problemphyxmini0716-roundstodisplayeddistance}
  \lean{PhyXMiniProblems.ProblemPhyXMini0716.RoundsToDisplayedDistance}
  \uses{def:physics:phyx-mini-0716:phyxminiproblems-problemphyxmini0716-lengthinmeters, def:physics:phyx-mini-0716:phyxminiproblems-problemphyxmini0716-binarystarsystem, def:physics:phyx-mini-0716:phyxminiproblems-problemphyxmini0716-answerchoice}
  The physical separation rounds to the displayed answer at its stated precision of one billion metres (half-width 5 * 10\textasciicircum{}8 m)
\end{definition}
\begin{proof}
  This is the declared model definition of Rounds To Displayed Distance.
\end{proof}
\begin{definition}[Is Unique Closest Answer]
  \label{def:physics:phyx-mini-0716:phyxminiproblems-problemphyxmini0716-isuniqueclosestanswer}
  \lean{PhyXMiniProblems.ProblemPhyXMini0716.IsUniqueClosestAnswer}
  \uses{def:physics:phyx-mini-0716:phyxminiproblems-problemphyxmini0716-lengthinmeters, def:physics:phyx-mini-0716:phyxminiproblems-problemphyxmini0716-binarystarsystem, def:physics:phyx-mini-0716:phyxminiproblems-problemphyxmini0716-answerchoice}
  The selected displayed distance is strictly nearer than every other choice
\end{definition}
\begin{proof}
  This is the declared model definition of Is Unique Closest Answer.
\end{proof}
% --- Archon declaration topology end ---
% --- Archon physics formalization source end ---
